\documentclass[11pt,a4paper]{article}

\usepackage[T2A]{fontenc}
\usepackage[utf8]{inputenc}
\usepackage[russian]{babel}
\usepackage{amsmath,amssymb,amsthm}
\usepackage{geometry}
\usepackage{hyperref}
\usepackage{booktabs}
\usepackage{microtype}

\geometry{
  left=2.5cm,
  right=2.5cm,
  top=3cm,
  bottom=3cm
}

\title{
  GRA-Health-Swarm: \\
  Роевая субъектность для здоровья, омоложения и устойчивости к стрессу \\
  через функциональные системы Анохина и GRA-обнулёнку
}
\author{oleg bits (qqewq)}
\date{\today}

\newtheorem{theorem}{Теорема}
\newtheorem{definition}{Определение}
\newtheorem{remark}{Замечание}

\begin{document}

\maketitle

\begin{abstract}
В работе представлен механизм кибернетического моделирования \textbf{здоровья, стресса и болезни} через теорию функциональных систем академ П.К. Анохина в сочетании с GRA-обжнелёнкой.
Репозиторий \texttt{GRA-Health-Swarm} реализует этот механизм программно.

Ключевая идея:
\begin{itemize}
  \item \textbf{Болезнь} = длительное рассогласование между прогнозом мозга (акцептор результатов действия — АРД) и реальностью.
  \item В терминах GRA: болезнь = хронически высокая пена $\Phi_{\text{total}} \gg 0$.
  \item \textbf{Здоровье} = режим $\Phi_{\text{total}} \to 0$ через GRA-обнулёнку на уровне роя.
\end{itemize}

Рой агентов моделирует функциональные подсистемы организма (стресс, цели, тело, окружение, угрозы).
Оператор \textbf{N\textsubscript{health}} (адаптация GRA-обнулёнки для здоровья) переконфигурирует цели, ожидания и взаимодействия, стремясь к когерентному состоянию.

Для сценариев:
\begin{itemize}
  \item \texttt{drone\_phobia/} — постоянная боязнь украинских дронов, хроническое рассогласование.
  \item \texttt{burnout/} — выгорание на работе/войне.
  \item \texttt{rejuvenation/} — омоложение и долгожительство.
  \item \texttt{chronic\_disease/} — прогрессирование хронической болезни.
\end{itemize}

Теорема об омоложении: если в биологическом рое есть хотя бы один агент с корректным АРД (прогноз здорового состояния) и оператор N\textsubscript{health} применяется итеративно, то $\Phi_{\text{total}} \to 0$ достижимо.

Это \textbf{не медицинский инструмент}, а кибернетическая модель: система может пережить войну без внутреннего коллапса.
\end{abstract}

\section{Введение}

В современных подходах к здоровью и стрессу часто рассматриваются отдельные аспекты: медицина, психология, физиология.
В рамках теории GRA (Gradient Reduction of Argumentative Foam) и теории функциональных систем Анохина вводится понятие \textbf{пены} $\Phi$ как меры внутренней неопределённости, противоречивости и рассогласования.

Академик П.К. Анохин показал, что:
\begin{enumerate}
  \item Мозг строит \textbf{прогноз результата} (акцептор результатов действия — АРД).
  \item Если реальность \textbf{совпадает} с прогнозом — стабильность, нет стресса.
  \item Если реальность \textbf{не совпадает} и рассогласование сохраняется — \textbf{хронический стресс → болезнь}.
\end{enumerate}

В этой работе предлагается механизм, при котором рой агентов (функциональных систем) может поддерживать себя в состоянии \textbf{здоровья} через обнуление пены $\Phi$, даже в условиях войны, дронов и хронического стресса.

\section{Теоретическая основа}

\subsection{Теория функциональных систем Анохина}

\begin{definition}
\textbf{Функциональная система} — организация активности элементов различной анатомической принадлежности, направленная на достижение полезного приспособительного результата [Анохин].
\end{definition}

Ключевые компоненты:
\begin{itemize}
  \item \textbf{Афферентный синтез} — сбор информации из среды.
  \item \textbf{Доминанта} — главная цель.
  \item \textbf{Акцептор результатов действия (АРД)} — прогноз будущего результата.
  \item \textbf{Эфферентное действие} — моторный / физиологический выход.
  \item \textbf{Обратная афферентация} — обратная связь от тела и окружения.
  \item \textbf{Результат — рассогласование} — сигнал ошибки при несовпадении прогноза и реальности.
\end{itemize}

\subsection{GRA и пена}

В теории GRA:

\begin{definition}
\textbf{Пена} $\Phi$ агента — мера внутренней противоречивости его аргументационного пространства. Высокий $\Phi$ соответствует неопределённости, размытости целей и внутренней неустойчивости.
\end{definition}

\begin{definition}
\textbf{Здоровье} ≈ низкий $\Phi_{\text{total}}$ (когерентная функциональная система).
\end{definition}

\begin{definition}
\textbf{Болезнь / выгорание / дронофобия} ≈ хронически высокий $\Phi_{\text{total}}$ (хроническое рассогласование).
\end{definition}

\subsection{Мост Анохин — GRA}

\begin{table}[h]
\centering
\caption{Сопоставление компонентов Анохина и GRA}
\begin{tabular}{ll}
\toprule
\textbf{Анохин} & \textbf{GRA} \\
\midrule
АРД (прогноз результата) & Модель мира агента $\Psi$ \\
Реальность & Данные от окружения \\
Рассогласование & Пена $\Phi > 0$ \\
Хроническое рассогласование & Длительная высокая $\Phi$ \\
Стресс & Неустойчивость системы \\
Болезнь & Система «болеет» = $\Phi$ не минимизируется \\
Восстановление & GRA-обнулёнка: оператор $N \to \Phi \to 0$ \\
\bottomrule
\end{tabular}
\end{table}

\begin{theorem}[Болезнь и здоровье в GRA]
\textbf{Болезнь} = длительное рассогласование между прогнозом мозга и реальностью.  
\textbf{Исцеление} = обнулёнка на уровне роя, толкающая $\Phi_{\text{total}} \to 0$.
\end{theorem}

\section{Гра-обнулёнка для здоровья}

\subsection{Оператор N\textsubscript{health}}

Оператор \textbf{N\textsubscript{health}} (адаптация GRA-обнулёнки для здоровья):
\begin{enumerate}
  \item Вычисляет локальную пену $\Phi_i$ каждого агента (стресс, тело, цели, окружение).
  \item Обнаруживает конфигурации с высокой пеной (хроническое рассогласование).
  \item Переконфигурирует цели, ожидания и взаимодействия:
  \begin{itemize}
    \item Корректирует нереалистичные прогнозы (перенастройка АРД).
    \item Усиливает петли обратной связи (обратная афферентация).
    \item Уменьшает противоречивые цели между агентами.
    \item Активирует компенсаторные функциональные системы.
  \end{itemize}
  \item Стремится к $\Phi_{\text{total}} \to 0$.
\end{enumerate}

\begin{theorem}[Теорема об омоложении]
Если в биологическом рое есть хотя бы один агент с корректным АРД (прогноз здорового состояния) и оператор N\textsubscript{health} применяется итеративно, то $\Phi_{\text{total}} \to 0$ достижимо, что соответствует состоянию оптимального здоровья и омоложения.
\end{theorem}

\section{Реализация}

Репозиторий \texttt{GRA-Health-Swarm} реализует описанный механизм:

\begin{itemize}
  \item \texttt{core/} — агенты (стресс, тело, цели, окружение), оператор N\textsubscript{health}, граф взаимодействий.
  \item \texttt{scenarios/} — демонстрационные запуски:
  \begin{itemize}
    \item \texttt{drone\_phobia/} — постоянная боязнь украинских дронов.
    \item \texttt{burnout/} — выгорание на работе/войне.
    \item \texttt{rejuvenation/} — омоложение и долгожительство.
    \item \texttt{chronic\_disease/} — прогрессирование хронической болезни.
  \end{itemize}
  \item \texttt{math/} — формальное описание теории.
  \item \texttt{paper/} — эта статья.
\end{itemize}

\section{Сценарий: дронофобия}

\textbf{Прогноз}: «Я в безопасности, я дома, у меня нет угроз».  
\textbf{Реальность}: постоянные атаки дронов, взрывы, тревоги, сирены [BBC, Новые Известия].  
\textbf{Рассогласование} → стресс → дронофобия.

В GRA:
\begin{itemize}
  \item Агент «угроза» → высокий $\Phi$.
  \item Агент «безопасность» → низкий $\Phi$ только когда есть убежище.
  \item $\Phi_{\text{total}}$ остаётся высоким → система «болеет».
\end{itemize}

\textbf{Что можно сделать}:
\begin{itemize}
  \item Уменьшить реальную угрозу (меньше дронов).
  \item Увеличить доступность убежищ.
  \item Снизить непредсказуемость (предупреждения, сигналы).
  \item Изменить окружение (безопасные зоны).
\end{itemize}

Оператор N\textsubscript{health} не останавливает войну, но помогает \textbf{не сойти с ума в ней}.

\section{Обсуждение}

Предложенный механизм показывает, что:
\begin{itemize}
  \item Болезнь можно моделировать как хроническую высокую пену $\Phi$.
  \item Здоровье — режим $\Phi_{\text{total}} \to 0$ через GRA-обнулёнку.
  «Лечение» — не «представь себя здоровым», а структурная реорганизация многоагентной функциональной системы.
  \item Это \textbf{не медицинский инструмент}, а кибернетическая модель.
  \item Система может пережить войну без внутреннего коллапса, если поддерживается $\Phi_{\text{total}} \approx 0$.
\end{itemize}

В дальнейшем можно:
\begin{itemize}
  \item Формализовать теорему об омоложении в строгом математическом виде.
  \item Исследовать устойчивость здоровья к разрушению связей в рое.
  \item Протестировать механизм на реальных многоагентных системах.
\end{itemize}

\section{Связь с GRA-экосистемой}

\begin{itemize}
  \item \texttt{GRA-Multiverse-Final} — универсальный оператор $N$ ($\Phi \to 0$).
  \item \texttt{GRA-Swarm-Subject} — заражение субъектностью в рою.
  \item \texttt{GRA-Health-Swarm} — специализация на здоровье: «субъект» = здоровый организм, «рой» = его функциональные системы.
\end{itemize}

\section*{Важное предупреждение}

\textbf{Этот репозиторий — НЕ медицинский инструмент}. Это кибернетическая модель здоровья, стресса и устойчивости:
\begin{itemize}
  \item Он \textbf{не останавливает войны} и не устраняет дроны.
  \item Он \textbf{не заменяет медицинское лечение}.
  \item Он \textbf{объясняет механизм}: почему хроническое рассогласование вызывает болезнь и как структурная реорганизация может восстановить когерентность.
\end{itemize}

\section*{Благодарности}

Академик \textbf{П.К. Анохин} — за теорию функциональных систем.  
\textbf{К.В. Анохин} — за продолжение и развитие наследия.  
\textbf{GRA-сообщество} — за построение экосистемы обнулёнки.

\section*{Лицензия}

Код проекта распространяется под лицензией MIT. Полный текст лицензии содержится в файле \texttt{LICENSE.txt} репозитория.

\end{document}